\documentclass{article}

% Required packages
\usepackage[utf8]{inputenc}
\usepackage[T1]{fontenc}
\usepackage{amsmath,amssymb,amsthm}
\usepackage{mathtools}
\usepackage{graphicx}
\usepackage{booktabs}
\usepackage{hyperref}
\usepackage{cleveref}
\usepackage{algorithm}
\usepackage{algpseudocode}
\usepackage{natbib}
\usepackage[margin=1in]{geometry}
\usepackage{tikz}
\usetikzlibrary{arrows.meta, positioning, shapes.geometric, calc}
\usepackage{xcolor}
\usepackage{enumitem}


% Theorem environments
\newtheorem{definition}{Definition}
\newtheorem{axiom}{Axiom}
\newtheorem{proposition}{Proposition}
\newtheorem{theorem}{Theorem}

\title{\textbf{Morph: Version Control When Pipelines Are Probabilistic}}

\author{%
  Raffi Krikorian\\
  Mozilla\\
  \texttt{raffi@mozilla.org}
}

\date{}

\begin{document}

\maketitle

\begin{abstract}
The tools we use to manage software were built for a world where code is
deterministic. Git tracks source files where identity is byte equality,
reproducibility means identical output, and merge is text reconciliation. None
of that holds when you are building with LLMs. Prompt outputs are stochastic.
Behavior depends on model version, retrieval corpus, and tool availability---not
just the files in the repo. And ``did it get better?'' is a statistical
question, not a diff.

\textsc{Morph} is a distributed version control system designed for this world.
It extends Git's content-addressed Merkle DAG with three additions:
pipelines (the sequences of prompt calls, tool invocations, retrieval steps,
and transforms that make up an LLM application), evaluation suites (versioned
definitions of what ``good'' means and how to measure it), and runs (permanent
execution receipts recording exactly what ran, in what environment, and what
it produced). A \textsc{Morph} commit bundles a pipeline with an evaluation
suite and scores. At merge time, \textsc{Morph} records the scores from both
parents and the scores the merged pipeline achieved---making it visible whether
the merge improved, held steady, or regressed on any metric.

The paper formalizes the pipeline model, defines what it means for one pipeline
version to improve on another, and works out the merge semantics that follow.
It describes a concrete v0 system, what the design looks like in practice for
agent-driven development and human-AI collaboration, and how the Actor
abstraction handles everything from a solo human to a swarm of agents using
the same DVCS machinery.
\end{abstract}

% ============================================================
\section{Introduction}
\label{sec:intro}

Software development is changing faster than the tools built to manage it.
Large language models are not autocomplete for code anymore---they are
components inside larger systems that retrieve context, call tools, generate
and modify code, run tests, and iterate~\citep{yao2023react,
jimenez2024swebench}. These pipelines---combining prompt calls, retrieval, tool
invocations, and deterministic transforms---are different from traditional source
code in ways that break the foundations of how we track changes.

Git~\citep{torvalds2005git}, the version control system that runs almost
everything, is built on three assumptions: (i)~identity is byte equality,
(ii)~reproducibility means identical output bytes, and (iii)~merge is syntactic
text reconciliation. All three break when the thing you are versioning is a
stochastic, environment-dependent transformation pipeline:

\begin{itemize}[nosep]
  \item \textbf{Same code, different outputs.} Run the same prompt twice and
    you get different answers. That is not a bug---it is how LLMs work. Git's
    assumption that identity means byte equality does not hold here.
  \item \textbf{Behavior depends on the environment.} Model version, sampling
    settings, retrieval corpus, and tool availability all shape what comes out.
    None of that is in the file tree.
  \item \textbf{``Did it get better?'' is a statistical question.} You cannot
    read a diff and know whether the pipeline improved. You have to run it,
    score it, and compare.
  \item \textbf{Agents produce patches, not just humans.} Which agent, which
    model, which prompt generated a change---and whether the result actually
    passes tests---matters for accountability. Git has no place to put any of
    that.
  \item \textbf{Merge can silently regress behavior.} Two branches can merge
    cleanly at the text level while the resulting pipeline performs worse than
    either parent.
\end{itemize}

\noindent This is not hypothetical. It is what happens every day in an
increasing share of software development, and it is why we need a different
kind of version control---one that records not just \emph{what the files are}
but \emph{who changed them}, \emph{what the pipeline did}, and \emph{whether it got better}.

\textsc{Morph} extends version control to \emph{recorded behavior}. It keeps
Git's core principles---every object identified by a hash of its contents,
organized in a Merkle DAG---while adding support for pipelines, evaluation
suites, and execution evidence. It is built for teams that have adopted the
improvement model: the shared assumption that each commit should leave the
pipeline at least as capable as before, the same way Git assumes you are
trying to write good code. Within that model, \textsc{Morph}'s job is to
record everything---every actor, every run, every score, every merge. A
\textsc{Morph} commit is not just a file snapshot; it is a record that a
specific pipeline, run against a specific evaluation suite, achieved specific
scores, with the execution receipts attached.

When a pipeline makes LLM calls, retrieves context, and runs tools, a merge
record cannot be purely syntactic. It has to include \emph{what the pipeline did}.
\textsc{Morph} records, at every merge, the scores both parents achieved and
the scores the merged pipeline achieved. What anyone does with that record is
up to them.

The contributions of this paper are:

\begin{enumerate}[nosep]
  \item A formal model of pipelines as composable, effect-aware computations
    supporting sequential chaining, parallel branching, and multi-actor
    attribution (\Cref{sec:programs}).
  \item A way to define what ``better'' means: evaluation suites as versioned
    contracts, scores as the result of running them, and an ordering that lets
    you compare any two pipeline versions (\Cref{sec:eval}).
  \item A merge record: \textsc{Morph} records the scores both parents achieved
    and the scores the merged pipeline achieved, with a formal property showing
    what it means when those scores improve (\Cref{sec:merge}).
  \item A concrete system design (\textsc{Morph} v0) built as a content-addressed
    object store with Git-style CLI tooling (\Cref{sec:system}).
\end{enumerate}

% ============================================================
\section{Background and Related Work}
\label{sec:related}

\textsc{Morph} builds directly on Git's architecture~\citep{torvalds2005git}: a
chain of commits where each one is identified by a hash of its contents,
organized into a Merkle DAG~\citep{merkle1988digital}, so history is
tamper-evident and cheap to compare. That design worked well for decades.
When code is deterministic, text is all you need to track. Prior work on
formalizing version control~\citep{roundy2005darcs,angiuli2014homotopical,swierstra2014semantics}
built more rigorous foundations for patch theory and merge semantics---but all
of it operates at the level of file content. \textsc{Morph} takes the same
impulse and applies it one layer up: to what the actors \emph{did} and what
the pipeline \emph{produced}, not just what the files \emph{say}.

ML tooling has partially recognized the gap. DVC~\citep{dvc2020} versions
datasets and pipelines alongside code; MLflow~\citep{zaharia2018mlflow} tracks
experiments and metrics across training runs. Evaluation frameworks like
HELM~\citep{liang2023helm} and SWE-bench~\citep{jimenez2024swebench} take
behavioral measurement seriously. But all of these treat evaluation as something
that happens \emph{around} version control, not inside it. Metrics get logged;
they do not become part of the commit record.

The closest thing practitioners have built is tools like Braintrust~\citep{braintrust2024},
LangSmith, and Humanloop---prompt versioning, evaluation pipelines, and CI hooks
that block a deploy when scores drop. That reflects a real need. But these are
centralized SaaS platforms that version individual prompts, not pipeline DAGs;
their merge blocking is a CI side-effect, not a native VCS operation; and they
produce no decentralized, verifiable evidence graph. They version prompts.
\textsc{Morph} records \emph{what pipelines did}---every actor, every run,
every score---and keeps the execution receipts to back it up.

LLM application frameworks like LangChain~\citep{chase2022langchain} and agent
paradigms like ReAct~\citep{yao2023react} define what the pipelines look like
and how agents operate~\citep{wang2024survey}. \textsc{Morph} is not a
competitor to any of them---it is the version control layer underneath them.

The pipeline formalism uses monadic effect theory~\citep{moggi1991notions,
wadler1995monads}---the same pattern behind \texttt{Promise} in JavaScript
and \texttt{Result} in Rust---to compose steps that carry side effects
without losing track of them. The reproducibility framing follows
Peng~\citep{peng2011reproducible}, updated for a setting where byte-identical
output is often impossible and re-running the same evaluation checks is the
achievable substitute.

% ============================================================
\section{Pipelines and How They Compose}
\label{sec:programs}

\subsection{What Gets Checked Out}

\textsc{Morph} models a workspace as three things held together:
\begin{equation}
  S = D \times C \times M,
\end{equation}
where $D$ is the document tree (your code, datasets, artifacts), $C$ is
execution context (intermediate results, caches), and $M$ is metadata (who
produced what, run identifiers). Think of it as the full snapshot of everything
in your working directory, plus the context needed to understand it. This is
the thing that gets checked out, modified, and committed.

\subsection{Actors}

Every worker---human, agent, or a human and agent working together---is an
\emph{Actor}:
\begin{equation}
  \text{Actor} = (\text{id},\; \text{type} \in \{\texttt{human},
  \texttt{agent}\},\; \text{env\_config} \cup \{\bot\}),
\end{equation}
where \texttt{env\_config} records the model, sampling settings, toolchain,
and runner for agent actors, and is $\bot$ for humans. The model does not
otherwise distinguish between human and agent workers. A human on a laptop, an
agent in a CI environment, and a human+agent pair in a Cursor session are all
just Actors with different types and environment configurations.

This matters: two agents working simultaneously in different environments is
exactly the same situation as two humans working on two laptops. Both check out
the tree, both work independently, both commit and push. That is the standard
distributed version control problem---and \textsc{Morph} inherits the standard
DVCS solution for it. No new theory needed. The \texttt{env\_config} on each
actor's nodes is just provenance---a record of what was running when---not a
coordination mechanism.

\subsection{What a Pipeline Is}

A pipeline takes a workspace state and produces a new one, plus a full record
of everything that happened along the way: which actor ran which step, which
model was called, what it returned, what tools ran, what failed and was retried.
Formally:
\begin{equation}
  P : S \to F(S),
\end{equation}
where $F$ is a wrapper that carries side effects (randomness, I/O, traces,
failures, retries) alongside the result, not just the final output.

Concretely, a pipeline is a DAG where each node is a typed operator and each
edge is either a data dependency (output of one node feeds into the next) or a
control dependency (ordering without data flow).

\begin{definition}[Pipeline]
A \emph{pipeline} $P = (V, \mathcal{E}, \kappa, \rho, \alpha, \varepsilon)$
consists of:
\begin{itemize}[nosep]
  \item $V$ --- a set of nodes
  \item $\mathcal{E} \subseteq V \times V$ --- directed edges forming a DAG
  \item $\kappa : V \to \{\texttt{prompt\_call}, \texttt{tool\_call},
    \texttt{retrieval}, \texttt{transform}, \texttt{identity},
    \texttt{review}\}$ --- node type, assigning each node its operator kind
  \item $\rho : V \to \mathcal{H} \cup \{\bot\}$ --- payload reference,
    pointing each node to its stored content (prompt template, tool schema,
    etc.) by hash, or $\bot$ if the node needs none
  \item $\alpha : V \to 2^{\mathcal{A}}$ --- attribution set, the set of Actor
    IDs that contributed to this node
  \item $\varepsilon : V \to \text{EnvConfig} \cup \{\bot\}$ --- per-node
    environment, recording which model and toolchain ran this node, or $\bot$
    for human-only nodes
\end{itemize}
\end{definition}

\noindent Throughout, $\bot$ denotes an absent or unattributed value. The same
symbol is reused deliberately wherever something is optional.

A word on \texttt{review} nodes. They represent an explicit acceptance or
modification decision---a human approving a diff in Cursor, or an agent
evaluating a candidate and choosing to accept, reject, or modify it. The
actor on a \texttt{review} node can be human, agent, or both. The node type
records the \emph{kind} of operation; the attribution set records \emph{who}
did it.

Attribution as a set handles the real-world mess. A human who accepts 80\% of
an agent's suggested diff and rewrites the other 20\% inline gets attribution
$\{\texttt{agent-1}, \texttt{human-1}\}$. One actor working alone gets a set
with one entry. A legacy node with no recorded attribution gets an empty
set---which is what $\bot$ reduces to.

\subsection{Composing Pipelines}

The tricky thing about chaining LLM steps is that each step does not just
return a value---it also produces side effects. A prompt call might fail and
retry. A tool invocation logs what it did. A retrieval step records latency.
Pipe outputs to inputs naively and you lose all of that bookkeeping.

A \emph{monad} solves this. You have almost certainly used one without calling
it that: \texttt{Promise} in JavaScript, \texttt{Optional} in Java,
\texttt{Result} in Rust. The pattern is the same everywhere---a wrapper that
carries extra context (errors, traces, retries) alongside the result, so that
context does not get dropped when you chain steps together.

For \textsc{Morph}, the monad gives us two operations:
\begin{align}
  \text{pure} &: A \to F(A), \\
  \text{bind} &: F(A) \to (A \to F(B)) \to F(B).
\end{align}

\noindent \texttt{pure} lifts a plain value into the pipeline so it can be
passed along without side effects. \texttt{bind} is how you chain steps: run
the first, take its output, pass it to the next. If you have used
\texttt{Promise.then()} in JavaScript or \texttt{flatMap} in Scala or Java,
that is bind---\textsc{Morph} just formalizes the same idea at the level of the
whole pipeline.

Sequential composition follows directly:
\begin{equation}
  (Q \circ P)(s) = \text{bind}(P(s),\; Q).
\end{equation}

\noindent The monad laws give you two things:
\begin{itemize}[nosep]
  \item \textbf{Associativity:} $(R \circ Q) \circ P = R \circ (Q \circ P)$
    --- you can refactor a pipeline without changing what it does.
  \item \textbf{Identity:} There exists a pipeline $I$ such that $I
    \circ P = P \circ I = P$ --- there is a meaningful no-op, which matters
    for defaults and incremental adoption.
\end{itemize}

For parallel branches---two actors working on the same input state, generating
candidates to be compared or merged---we fan out, run both, and reconcile:
\begin{equation}
  \text{branch}(P, Q) = J \circ (P \otimes Q) \circ \Delta,
\end{equation}
where $\Delta(s) = (s, s)$ duplicates the starting state, $P \otimes Q$ runs
both branches independently (this is \texttt{Promise.all()}), and $J$ is an
explicit join node that decides what to do with both results---pick the one with
better test coverage, merge the best parts, or hand off to a \texttt{review}
node. The join step is itself a node in the pipeline with its own attribution
and record of who ran it.

\subsection{Multiple Actors, Concurrent Work}

The concurrent case---two humans on two laptops, two agents in two cloud
environments, a human and an agent each on their own branch---is just DVCS.
Each actor checks out the tree, works independently, commits, and pushes.
\textsc{Morph} records the scores from each branch and what the merge achieved.
The standard distributed version control model handles everything else.

The interesting cases are within a single session, where multiple actors
contribute to one pipeline:

\begin{itemize}[nosep]
  \item \textbf{Sequential handoff:} Actor A does retrieval, actor B writes the
    code, actor C runs verification. This is just sequential composition
    ($C \circ B \circ A$) with different actors on different nodes.
    Attribution records who did each step; \texttt{env\_config} records what
    environment each ran in.
  \item \textbf{Parallel candidates:} Two actors generate independent candidate
    patches from the same starting state. A \texttt{review} node picks or
    merges them. This is the branch composition above---the standard Cursor
    multi-composer workflow.
  \item \textbf{Human + agent on one node:} A human accepts and edits an
    agent's diff. This is a \texttt{review} node with attribution
    $\{\texttt{agent}, \texttt{human}\}$. No special case needed.
\end{itemize}

\noindent Same record in all three cases: who contributed what, in what
environment, producing what output. The scores are a property of the composed
result, not any individual step.

One thing \textsc{Morph} cannot do from the record alone is decompose
blame to individual actors. If the pipeline scores poorly,
\textsc{Morph} knows the score and knows who touched which nodes, but it cannot
automatically tell you which actor caused it. Contributions are
entangled---agent B's code change may only make sense given what
agent A retrieved. Figuring this out means running the pipeline again
with one actor's contribution removed. That is future work.

% ============================================================
\section{Defining ``Better'': Evaluation and Scores}
\label{sec:eval}

\subsection{Evaluation Suites}

In practice, an evaluation suite is the answer to: what does ``good'' mean for
this pipeline? In a Claude Code workflow, it might be: does the generated patch
pass the test suite? Does it stay under 500ms p95 latency? Does it call any
forbidden APIs? In a Cursor session building a RAG pipeline, it might be:
does retrieval hit the right documents in the top-3? Does the generated answer
score above 0.85 on faithfulness? Does the whole thing cost under \$0.02 per
query?

In code, an evaluation suite $T$ takes a workspace state and produces
observations:
\begin{equation}
  \text{Eval}_{T,E} : S \to F(\text{Obs}_T),
\end{equation}
where $\text{Obs}_T$ includes raw metric samples, test execution logs, timing
data, and anything else needed to produce scores. Running an evaluation has
side effects too---compiling code, calling a judge model, collecting human
ratings---so it lives in the same $F$ wrapper as pipelines. Each metric
$m \in T$ specifies a scoring function, a direction (higher is better, or
lower is better), and how to aggregate samples into a final score (mean, lower
confidence bound, percentile threshold, pass/fail).

\textsc{Morph} supports two kinds of evaluation:

\begin{itemize}[nosep]
  \item \textbf{Artifact evaluation:} Does the output do what it should?
    Unit tests pass, integration tests pass, static analysis clean, benchmarks
    within bounds. This is the \texttt{pytest} or \texttt{cargo test} you are
    already running---\textsc{Morph} just makes the results a versioned
    object you can compare across commits.
  \item \textbf{Process evaluation:} Did the pipeline behave acceptably while
    getting there? Cost under budget, latency within bounds, no forbidden tool
    calls, no policy violations. Things that do not show up in the output but
    matter for deployment.
\end{itemize}

\noindent Both matter. A patch that passes tests but costs \$10 per request
is not a good merge, even if the text diff is clean.

There is one trap worth calling out. An agent that controls its own test files
could modify them to manufacture a perfect score. Each evaluation suite
declares a \emph{fixture source} to prevent this: where the test data comes
from. Options are \texttt{candidate-sourced} (the produced tree),
\texttt{base-sourced} (a specified baseline), \texttt{pinned} (locked to a
specific hash), or \texttt{external}. Without this, a capable enough agent
could just rewrite the tests.

\subsection{Scores}

For an evaluation suite $T$ with metrics $\{m_1, \ldots, m_k\}$, a scorecard
is one number per metric---accuracy, latency, cost, whatever $T$ defines. We
write $V_T$ for the set of all possible scorecards under $T$:
\begin{equation}
  V_T = \prod_{m \in T} V_m.
\end{equation}

One scorecard is at least as good as another if it scores better on
\emph{every} metric (e.g., higher for accuracy, lower for latency):
\begin{equation}
  x \leq_T y \iff \forall m \in T,\; x_m \leq_m y_m.
\end{equation}

Running a pipeline $P$ from starting state $s_0$ against suite $T$ produces a
scorecard. Because a pipeline's outputs vary across runs, a single execution
gives you a sample, not a ground truth. The aggregation method---mean over $n$
runs, lower confidence bound at level $\alpha$, p95 latency, pass/fail---is
part of the evaluation suite definition, fixed and hashed like everything else.
$\text{cert}_T(P, s_0)$ denotes the resulting aggregate scorecard:
\begin{equation}
  \text{cert}_T(P, s_0) \in V_T.
\end{equation}

A score difference between two pipelines might be a real improvement or
sampling noise. The formalism treats $\text{cert}_T$ as the declared aggregate
and leaves the question of significance to the evaluation suite designer; the
discussion covers what that means in practice.

\begin{definition}[Improvement Ordering]
Pipeline $Q$ \emph{improves on} pipeline $P$ under suite $T$ and starting
state $s_0$, written $P \preceq_{T,s_0} Q$, if its scorecard is at least as
good on every metric:
\begin{equation}
  \text{cert}_T(P, s_0) \leq_T \text{cert}_T(Q, s_0).
\end{equation}
\end{definition}

\begin{proposition}
The improvement ordering $\preceq_{T,s_0}$ is a preorder: a pipeline always
improves on itself (reflexivity), and if $Q$ improves on $P$ and $R$ improves
on $Q$, then $R$ improves on $P$ (transitivity).
\end{proposition}
\begin{proof}
Reflexivity follows from reflexivity of $\leq_m$ for each metric. Transitivity
follows from transitivity of $\leq_m$ applied componentwise.
\end{proof}

\noindent This turns ``did it get better?'' from a judgment call into a
comparison with a well-defined answer.

% ============================================================
\section{Commits and Merge}
\label{sec:merge}

\subsection{What's in a Commit}

A Git commit says: here is what the files looked like. A \textsc{Morph} commit
says: here is a pipeline, the evaluation suite it was run against, and the
scores it got. The execution receipts are attached.

\begin{definition}[Commit]
A \textsc{Morph} commit is a tuple
\begin{equation}
  c = (\text{tree\_hash},\; \text{program\_id},\; T,\; v,\; \text{parents},\;
  \text{env\_constraints},\; \text{evidence\_refs}),
\end{equation}
where $\text{tree\_hash}$ is the root hash of the working directory tree (same
role as Git's tree object), $\text{program\_id}$ is the content hash of the
pipeline DAG, $T$ is an evaluation suite (by hash), $v \in V_T$ is the scores
the pipeline achieved, $\text{parents}$ are parent commit hashes forming the
Merkle DAG, $\text{env\_constraints}$ records the environment in which the
scores were captured, and $\text{evidence\_refs}$ are hashes of the
supporting run and trace objects.
\end{definition}

The scores $v$ are not decorative. They are what the pipeline produced when it
was committed, and the reference point for any future merge. When the pipeline
and evaluation suite are unspecified, they default to the identity pipeline and
empty suite---so \textsc{Morph} degrades gracefully to a plain VCS. Same
tamper-evident history as Git. Just more in each commit.

\subsection{How Merge Works}

Git merge records a text reconciliation. \textsc{Morph} is built on the
assumption that you are merging because you believe the result is better than
either parent. The merge record captures everything needed to verify that
assumption---or to see exactly where it fell short:

\begin{enumerate}[nosep]
  \item \textbf{Structural merge} of pipeline DAGs (analogous to three-way
    text merge).
  \item \textbf{Combine evaluation suites.} If parents have suites $T_1$
    and $T_2$, the merge suite is $T = T_1 \uplus T_2$---the full set of
    metrics from both, kept separate by metric ID so definitions cannot
    silently collide.
  \item \textbf{Record the bar.} Embed each parent's scores into $V_T$
    and record the best from either parent on every metric:
  \begin{equation}
    v_{\text{ref}} = \text{embed}(v_1) \sqcup \text{embed}(v_2).
  \end{equation}
  The embedding $\text{embed} : V_{T_i} \to V_T$ maps each metric's value
  from the parent suite to the corresponding component in the merged suite,
  and assigns $\bot_m$ (the worst possible value for metric $m$) to any
  metric in $T \setminus T_i$. This requires each $V_m$ to have a least
  element $\bot_m$, which we take as an axiom: for a minimize-direction
  metric, $\bot_m = +\infty$; for a maximize-direction metric, $\bot_m = -\infty$.
  The operator $\sqcup$ then takes the componentwise maximum under each
  metric's direction.
  \item \textbf{Record improvement.} Run the merged pipeline $R$ and record
    its scores relative to $v_{\text{ref}}$:
  \begin{equation}
    \text{cert}_T(R, s_0) \text{ vs. } v_{\text{ref}}.
    \label{eq:merge-dom}
  \end{equation}
\end{enumerate}

\begin{theorem}[What a Good Merge Record Shows]
\label{thm:monotone}
If a merge commit records $\text{cert}_T(R, s_0) \geq_T v_{\text{ref}}$,
then the merged pipeline $R$ scores at least as well as both parents on every
metric each parent measured. For parent commits $c_1 = (\ldots, T_1, v_1, \ldots)$ and $c_2 =
(\ldots, T_2, v_2, \ldots)$:
\begin{equation}
  \text{cert}_{T_1}(R, s_0) \geq_{T_1} v_1 \quad \text{and} \quad
  \text{cert}_{T_2}(R, s_0) \geq_{T_2} v_2.
\end{equation}
\end{theorem}

\begin{proof}
By construction, $\text{embed}(v_i)$ agrees with $v_i$ on every metric in
$T_i$ and assigns $\bot_m$ elsewhere. The join $v_{\text{ref}} =
\text{embed}(v_1) \sqcup \text{embed}(v_2)$ therefore satisfies
$v_{\text{ref}} \geq_{T_i} \text{embed}(v_i) \geq_{T_i} v_i$ for $i \in
\{1, 2\}$ (where $\geq_{T_i}$ denotes restriction to the metrics of $T_i$).

If $\text{cert}_T(R, s_0) \geq_T v_{\text{ref}}$, then in particular for
every metric $m \in T_1$:
\[
  \text{cert}_T(R, s_0)_m \geq_m (v_{\text{ref}})_m \geq_m (v_1)_m,
\]
which gives $\text{cert}_{T_1}(R, s_0) \geq_{T_1} v_1$, and symmetrically
for $T_2$.
\end{proof}

If you look at a merge record and the merged pipeline scored at or above both
parents, nothing was lost. If the record shows a drop on any metric, you know
exactly where and by how much. Either way, the record is there.

\subsection{Metric Retirement}

The merge record compares against the parent suites. That works when the
parent metrics still make sense. Usually they do. But sometimes a branch
changes what the pipeline fundamentally does---switching retrieval strategy,
replacing a model, restructuring the task entirely. The old metrics may no
longer be the right thing to measure. Comparing the new pipeline against them
would be like measuring a search engine's latency with a benchmark designed for
a chatbot: the number is real, the comparison is meaningless.

\textsc{Morph} handles this with \emph{metric retirement}. A commit can declare that one or more metrics from a parent's evaluation suite
no longer apply, with a stated reason. A retired set $D \subseteq T_1 \uplus
T_2$ is removed from the merge suite before recording the reference scores:
\begin{equation}
  T = (T_1 \uplus T_2) \setminus D, \quad
  v_{\text{ref}} = \text{embed}(v_1) \sqcup \text{embed}(v_2)
    \text{ projected onto } T.
\end{equation}

The retirement is explicit and attributed: any merge that retires metrics must
include a \texttt{review} node. The review actor---human, agent, or both---is
recorded as the one who made the call. This is not a human requirement; it is
an attribution requirement. The \texttt{review} node records who decided, in
what environment, and why. That record persists like everything else.

The merge record still shows improvement on every metric that survived.
\textsc{Morph} does not prevent someone from retiring every metric---that is a
policy decision, not a recording one. What it does is make every retirement
visible, attributed, and permanent in the object graph.

\section{How \textsc{Morph} Works}
\label{sec:system}

\subsection{How Objects Are Stored}

Like Git, \textsc{Morph} stores everything as content-addressed, permanent
nodes in a Merkle DAG using SHA-256---meaning every object is identified by
a hash of its contents, so the history cannot be tampered with. The object
types extend Git's \texttt{blob}/\texttt{tree}/\texttt{commit} with new
additions:

\begin{table}[t]
\centering
\caption{Object types in \textsc{Morph} compared to Git. Items marked with
$\star$ are new to \textsc{Morph}; others generalize Git equivalents.}
\label{tab:objects}
\begin{tabular}{@{}llp{6.5cm}@{}}
\toprule
\textbf{Object} & \textbf{Git Analog} & \textbf{Description} \\
\midrule
Blob & blob & Atomic content: prompt templates, tool schemas, configs. Open
  \texttt{kind} field. \\
Tree & tree & Structured grouping of objects (directory analog). \\
Pipeline$^\star$ & --- & DAG of typed operators with data/control edges. \\
EvalSuite$^\star$ & --- & Versioned evaluation definition: cases, metrics,
  thresholds, fixture sources. \\
Commit & commit & Tree hash $+$ pipeline hash $+$ eval suite $+$ scores
  $+$ evidence refs. \\
Run$^\star$ & --- & Execution receipt: environment, inputs, outputs, metrics,
  trace ref, actor identity. \\
Trace$^\star$ & --- & Sequence of typed, addressable events within a run. \\
Annotation$^\star$ & --- & Metadata on any object without changing its hash. \\
\bottomrule
\end{tabular}
\end{table}

\subsection{Runs and Traces: Keeping the Receipt}

When Claude Code applies a patch, something happened. A specific model ran, with
specific sampling settings, on a specific input state, and produced a specific
output. Tests either passed or failed. Latency was recorded. Maybe a tool call
was made mid-generation. Right now, most of that disappears the moment the
session ends. \textsc{Morph} keeps it.

A \textbf{Run} is a permanent receipt of a single pipeline execution. It
records: the pipeline hash (so you know exactly what ran), the full environment
(model identifier and version, sampling settings, toolchain), the input state,
the output artifacts, the observed scores, a trace reference, and the actor's
identity. Everything is there. If you need to re-run the evaluation or dig
into why a merge looked the way it did, the Run has it.

A \textbf{Trace} is the detailed record inside a Run---a sequence of typed
events, each with a unique identifier. A prompt call event records the rendered
prompt and the raw completion. A tool call event records the invocation and
result. A retry records what failed and why. Individual events within a trace
can be annotated independently---so a reviewer can flag a specific tool call or
a specific model output without touching anything else.

Runs never modify commits. They are evidence that accumulates over time. This
means you can run a pipeline ten times and have ten Runs attached to the same
commit---each one adding to the picture of how it behaves, without rewriting
history.

\subsection{Annotations}

\textbf{Annotations} attach metadata to any object (or a specific event within
a trace) without altering its hash. They enable human ratings, bookmarks, tags,
cross-references, and attribution overlays. Annotations are the extensibility
mechanism: higher-level tools can layer rich metadata onto the permanent object
graph without requiring changes to core object types.

\subsection{Separation of Concerns: \textsc{Morph} Does Not Execute}

\textsc{Morph} does not execute pipelines or run evaluations. All LLM calls,
tool execution, and test runs happen in external tools---Cursor, Claude Code,
your CI system. \textsc{Morph} just stores what they report: permanent
content-addressed objects, execution receipts, scores, and the record of what
happened at every merge. Commands like \texttt{morph run record} and
\texttt{morph eval record} \emph{ingest} results; they do not run anything.
The VCS layer stays thin and out of the way.

\subsection{Integration: Editors, Agents, and the Filesystem}

\textsc{Morph} is designed to be integrated into the tools where AI-assisted
development actually happens---editors like Cursor and agentic coding tools like
Claude Code---not run alongside them as a separate process.

\paragraph{The editor path.}
In a VS Code-based environment like Cursor, the integration is a VS Code
extension. VS Code's document model already serializes all writes to open
files---when a human is typing and an agent is applying a diff simultaneously,
the editor queues them. There is no raw filesystem collision at the content
level. \textsc{Morph} hooks into VS Code's document change events, which fire
on every modification with the diff and a flag indicating whether the change
came from a user action or a programmatic edit (agents use
\texttt{workspace.applyEdit()}). The extension reads that flag, tags attribution
accordingly, and records each change as a trace event. The pipeline DAG builds
incrementally as events arrive. A commit is just: snapshot the current tree,
hash it, seal the DAG assembled so far.

\paragraph{The filesystem path.}
When agents write directly to the filesystem---headless Claude Code sessions,
subprocess invocations, scripts running outside the editor---the document
model's serialization is gone. Two agents could write to the same file
simultaneously.

Well-behaved agents already handle this. Claude Code writes a
file, reads it back, compiles, runs tests, and checks whether the world is in
the state it expected. If another agent wrote something conflicting in between,
the compile or test step catches it. The agent does not need to know about the
concurrent write---it just sees a broken state and responds to that. The
write-read-verify loop is the conflict detection mechanism.

\textsc{Morph}'s role is not to intercept or serialize filesystem access.
It is to record what the actor declares when it commits---which pipeline ran,
what the tree hash is, and what the evaluation results were. If two agents
collided and left the tree in a broken state, the evaluation scores reflect
that. The record shows exactly what happened and when.

Agents are responsible for verifying their own writes before committing. That
is already the convention. \textsc{Morph} records whether they did.



\subsection{CLI Design}

The CLI mirrors Git where possible:

\smallskip
\begin{center}
\begin{tabular}{@{}ll@{}}
\texttt{morph init} & Initialize repository \\
\texttt{morph add .} & Stage working-space changes \\
\texttt{morph commit -m "msg"} & Commit with eval contract \\
\texttt{morph merge <branch>} & Merge with score recording \\
\texttt{morph run record <file>} & Ingest execution receipt \\
\texttt{morph eval record <file>} & Ingest evaluation results \\
\texttt{morph annotate <hash>} & Attach metadata \\
\texttt{morph log} / \texttt{status} & Inspect history and state \\
\end{tabular}
\end{center}

% ============================================================
\section{Axioms}
\label{sec:axioms}

\textsc{Morph} rests on a small set of axioms. These are not philosophical
claims. They are the design constraints the rest of the system depends on.
Violate any of them in an implementation and something breaks.

\begin{axiom}[Content-Addressed, Immutable Objects]
Every object is identified by a cryptographic hash of its contents, the same
way Git works. If the hash matches, the content is what you think it is.
History cannot be tampered with.
\end{axiom}

\begin{axiom}[Evidence Does Not Rewrite History]
Runs and evaluation results never mutate prior commits. New evidence produces
new objects; old commits stay exactly as they were. You can always go back and
see what a commit recorded at the time it was made.
\end{axiom}

\begin{axiom}[Pipeline Steps Compose Cleanly]
There is a well-defined way to chain pipeline steps sequentially and run them
in parallel---analogous to \texttt{Promise.then()} for chaining and
\texttt{Promise.all()} for concurrency---such that the results stay consistent
when you refactor or restructure. You can reorganize a pipeline without
changing what it does, and there is a meaningful no-op step.
\end{axiom}

\begin{axiom}[Evaluation Suites are Explicit Contracts]
``Better'' is not implicit. Every evaluation suite declares its metrics,
their directions, their fixture sources, and how samples get aggregated into
scores. All of it is versioned and hashed like everything else.
\end{axiom}

\begin{axiom}[Scores are Partially Ordered]
A scorecard is one number per metric. One scorecard dominates another only if
it wins on every metric simultaneously. If pipeline A is more accurate but
slower than pipeline B, they are incomparable---neither dominates. You only
get a clean winner when one pipeline is at least as good on everything.
\end{axiom}

\begin{axiom}[Merge Records Scores From Both Parents]
A merge commit records what both parents achieved and what the merged pipeline
achieved. This is the record that lets anyone verify, for any merge in history,
whether the shared assumption held.
\end{axiom}

\begin{axiom}[Environment is Part of the Record]
Every run records the environment that produced its scores---model version,
sampling settings, toolchain. Without this, scores from different environments
are not comparable. A pipeline running against \texttt{gpt-4o} and the same
pipeline running against a fine-tuned local model are different things, even
if the code is identical.
\end{axiom}

\begin{axiom}[Reproducibility Means Re-Running the Checks]
You cannot get byte-identical outputs from an LLM pipeline. Reproducibility
means something weaker: re-running the same evaluation suite in the same
declared environment produces aggregate scores consistent with the recorded
value, within the expected variance of the aggregation method chosen.
\end{axiom}

% ============================================================
\section{Discussion}
\label{sec:discussion}

\paragraph{This is what CI is already trying to do.}
CI pipelines already run tests on proposed changes. GitHub Actions blocks a
merge if tests fail. Braintrust blocks a deploy if eval scores drop. These
tools share the same underlying assumption \textsc{Morph} is built on: that
you are trying to get better. The difference is that \textsc{Morph} records
the contract, the evidence, and the outcome as permanent versioned objects,
not ephemeral pipeline logs. Failed CI runs evaporate; \textsc{Morph} runs
persist as evidence you can audit months later. Any enforcement sits on top.
\textsc{Morph} just makes sure the record is always there.

\paragraph{Agent accountability, finally.}
When Claude Code applies a patch today, there is no clean record of which
model generated it, what prompt produced it, what it cost, or whether it
passed any checks. Something breaks in production. You go digging through
session logs hoping something was captured.

\textsc{Morph} changes that. Every committed patch has a Run attached: the
agent identity, the model version, the full execution trace, and the scores.
That chain---from agent action to recorded outcome---does not currently exist
anywhere in agent-driven development.

\paragraph{The stats are minimal in v0.}
The current implementation uses mean aggregation, worst-case scores, percentile
thresholds, and lower confidence bounds. That is enough to be useful. Richer
methods---Bayesian comparison, two-sample equivalence testing, human-in-the-loop
evaluation---all fit the evaluation interface and are deferred to future work.
The interface is clean: swapping in a better scoring method does not require
changing the commit format or the merge procedure.

\paragraph{Built to distribute.}
Because all objects are content-addressed and permanent, distributed settings
are just Git's push/pull model applied to a richer object store. The v0
implementation is local-only; a remote protocol and federated verification
are planned.

\paragraph{On Git compatibility.}
An implementation of \textsc{Morph} can use Git's object store directly.
The content-addressed Merkle DAG that Git uses is the right substrate for
\textsc{Morph}'s objects too---blobs, trees, and commits map cleanly, and
the new object types (pipelines, eval suites, runs, traces) are just additional
objects in the same store. \textsc{Morph} and Git can share a repository:
Git tracks source files, \textsc{Morph} tracks pipelines and score history,
same underlying object store. Adoption is incremental. Leave the pipeline and
evaluation suite fields empty in a \textsc{Morph} commit and it degenerates
into a standard tree-hash commit.

\paragraph{Scores are statistics, not facts.}
Every score is an aggregate over runs that vary---a mean, a confidence bound,
a percentile. An observed improvement of 0.02 accuracy might be real or it
might be noise, depending on sample count and pipeline variance. \textsc{Morph}
records the aggregate and the method that produced it. Teams that need
stronger guarantees should use lower confidence bounds or equivalence tests
rather than raw means; the record makes the choice visible either way.

\paragraph{The credit assignment problem.}
When multiple actors contribute to a single pipeline, \textsc{Morph} can tell
you who touched which nodes, but it cannot automatically tell you which actor
caused a score change. The scores are a property of the composed pipeline, not
individual contributions. Agent A's retrieval step shapes what agent B can
generate; agent B's output shapes what agent C verifies. These contributions
are entangled, not independent.

The standard approach---Shapley-value decomposition---requires running the
pipeline counterfactually with each actor's contribution removed. For LLM
pipelines, that is ill-defined: removing agent A's retrieval does not produce
a meaningful baseline because agent B's generation was conditioned on that
retrieval. You cannot swap out one actor and expect the rest of the pipeline
to behave as if it had not happened. For actors on separate branches this is
not a limitation---each branch has its own scores. For tightly coupled
cooperative work, the evaluation suite has to test the composed result as a
whole, and credit decomposition remains an open research problem.

\paragraph{Known limitations.}
Three open problems worth naming.

First, \textsc{Morph} assumes evaluation suites are honest---that the metrics
actually measure what you think they measure. A bad eval produces meaningless
scores. Garbage in, garbage out; \textsc{Morph} makes it auditable garbage,
but it is still garbage.

Second, structural merge of pipeline DAGs (step 1 of the merge procedure) is
under-specified. Three-way text merge has decades of tooling behind it.
Pipeline-graph merge does not. Even defining what a correct pipeline merge
means---when two branches both modified the same prompt node in different
directions---is an open problem. This is the hardest gap in \textsc{Morph} v0.

Third, a recorded score difference may be real improvement or sampling noise.
The formalism treats $\text{cert}_T$ as the declared aggregate and leaves
significance to the evaluation suite designer. A future version could expose
a significance threshold as a first-class field in the suite, making the
required confidence level part of the versioned contract rather than buried
in the aggregation method.

\paragraph{Distribution shift.}
Evaluation suites are pinned by hash, which is the right design for
auditability. But the world changes: retrieval corpora age, model providers
update APIs silently, user query distributions drift. A high score recorded
six months ago against a pinned suite may say nothing about current
performance. \textsc{Morph} does not solve this---no version control system
can. What it does is make the evaluation suite explicit and versioned, so
when a team updates their suite to reflect a shifted distribution, the change
is visible in the history with full attribution. The gap between a pinned
historical eval and a current one is a deliberate design choice that teams
can see, not a silent assumption they cannot.

\paragraph{Metric retirement and policy.}
Metric retirement is a necessary escape valve but also a potential hole.
Nothing in the formal model prevents an actor from retiring every metric.
That is intentional. \textsc{Morph}'s job is not to block things---it is to
record everything. Every retirement is a \texttt{review} node in the history:
who made the call, in what environment, and why. Organizations that want to
restrict who can retire which metrics can build that on top, the same way
code review policies sit on top of Git without being part of Git.

% ============================================================
\section{Conclusion}
\label{sec:conclusion}

\textsc{Morph} is built on a simple assumption: that teams using it are trying
to improve. Each commit should leave the pipeline at least as capable as before.
Each merge should produce something better than either parent. That is not a
constraint \textsc{Morph} enforces---it is the model actors buy into when they
use it, the same way Git assumes you are trying to write good code.

Within that model, \textsc{Morph}'s job is to record everything. It extends
Git's content-addressed Merkle DAG to capture what actors actually do: the
pipelines they run, the evaluation suites they run them against, and the scores
that come back. Every merge records what both parents achieved and what the
merged pipeline achieved. The Actor abstraction puts humans, agents, and
human-agent pairs on equal footing---same DVCS model, same commit mechanics,
same merge record. In a Cursor session, the pipeline DAG builds incrementally
from every change event. A commit is just a named snapshot of that stream.

When the software you are building makes LLM calls, retrieves context, and
runs tools---and the same code produces different outputs depending on which
model is running---``what changed'' is not just bytes. It is who did what,
under what conditions, and what the pipeline produced. \textsc{Morph} records
all of it.

\bibliographystyle{plainnat}

\begin{thebibliography}{20}

\bibitem[Angiuli et~al.(2014)]{angiuli2014homotopical}
Carlo Angiuli, Edward Morehouse, Daniel~R. Licata, and Robert Harper.
\newblock Homotopical patch theory.
\newblock In \emph{Proceedings of the 19th ACM SIGPLAN International Conference
  on Functional Programming (ICFP)}, pages 243--256, 2014.

\bibitem[Braintrust(2024)]{braintrust2024}
Braintrust.
\newblock {Braintrust: The AI observability platform}.
\newblock \url{https://www.braintrust.dev}, 2024.

\bibitem[Chase(2022)]{chase2022langchain}
Harrison Chase.
\newblock {LangChain}.
\newblock \url{https://github.com/langchain-ai/langchain}, 2022.

\bibitem[Chiang et~al.(2024)]{chiang2024chatbotarena}
Wei-Lin Chiang, Lianmin Zheng, Ying Sheng, Anastasios~Nikolas Angelopoulos,
  Tianle Li, Dacheng Li, Hao Zhang, Banghua Zhu, Michael Jordan,
  Joseph~E. Gonzalez, and Ion Stoica.
\newblock Chatbot {Arena}: An open platform for evaluating {LLMs} by human
  preference.
\newblock In \emph{Proceedings of the 41st International Conference on Machine
  Learning (ICML)}, 2024.

\bibitem[DVC(2020)]{dvc2020}
{DVC Team}.
\newblock {DVC}: Data version control.
\newblock \url{https://dvc.org}, 2020.

\bibitem[Jimenez et~al.(2024)]{jimenez2024swebench}
Carlos~E. Jimenez, John Yang, Alexander Wettig, Shunyu Yao, Kexin Pei, Ofir
  Press, and Karthik Narasimhan.
\newblock {SWE-bench}: Can language models resolve real-world {GitHub} issues?
\newblock In \emph{The Twelfth International Conference on Learning
  Representations (ICLR)}, 2024.

\bibitem[Liang et~al.(2023)]{liang2023helm}
Percy Liang, Rishi Bommasani, Tony Lee, Dimitris Tsipras, Dilara Soylu,
  Michihiro Yasunaga, Yian Zhang, Deepak Narayanan, Yuhuai Wu, Ananya Kumar,
  et~al.
\newblock Holistic evaluation of language models.
\newblock \emph{Transactions on Machine Learning Research}, 2023.

\bibitem[Merkle(1988)]{merkle1988digital}
Ralph~C. Merkle.
\newblock A digital signature based on a conventional encryption function.
\newblock In \emph{Advances in Cryptology---CRYPTO '87}, pages 369--378.
  Springer, 1988.

\bibitem[Moggi(1991)]{moggi1991notions}
Eugenio Moggi.
\newblock Notions of computation and monads.
\newblock \emph{Information and Computation}, 93(1):55--92, 1991.

\bibitem[Peng(2011)]{peng2011reproducible}
Roger~D. Peng.
\newblock Reproducible research in computational science.
\newblock \emph{Science}, 334(6060):1226--1227, 2011.

\bibitem[Roundy(2005)]{roundy2005darcs}
David Roundy.
\newblock Darcs: Distributed version management in {Haskell}.
\newblock In \emph{Proceedings of the ACM SIGPLAN Workshop on Haskell}, pages
  1--4, 2005.

\bibitem[Swierstra and L\"{o}h(2014)]{swierstra2014semantics}
Wouter Swierstra and Andres L\"{o}h.
\newblock The semantics of version control.
\newblock In \emph{Proceedings of the 2014 ACM International Symposium on New
  Ideas, New Paradigms, and Reflections on Programming \& Software (Onward!)},
  pages 43--54, 2014.

\bibitem[Torvalds(2005)]{torvalds2005git}
Linus Torvalds.
\newblock {Git}: Fast version control system.
\newblock \url{https://git-scm.com}, 2005.

\bibitem[Wadler(1995)]{wadler1995monads}
Philip Wadler.
\newblock Monads for functional programming.
\newblock In \emph{Advanced Functional Programming}, volume 925 of
  \emph{Lecture Notes in Computer Science}, pages 24--52. Springer, 1995.

\bibitem[Wang et~al.(2024)]{wang2024survey}
Lei Wang, Chen Ma, Xueyang Feng, Zeyu Zhang, Hao Yang, Jingsen Zhang, Zhiyuan
  Chen, Jiakai Tang, Xu~Chen, Yankai Lin, Wayne~Xin Zhao, Zhewei Wei, and
  Ji-Rong Wen.
\newblock A survey on large language model based autonomous agents.
\newblock \emph{Frontiers of Computer Science}, 18(6):186345, 2024.

\bibitem[Yao et~al.(2023)]{yao2023react}
Shunyu Yao, Jeffrey Zhao, Dian Yu, Nan Du, Izhak Shafran, Karthik Narasimhan,
  and Yuan Cao.
\newblock {ReAct}: Synergizing reasoning and acting in language models.
\newblock In \emph{International Conference on Learning Representations
  (ICLR)}, 2023.

\bibitem[Zaharia et~al.(2018)]{zaharia2018mlflow}
Matei Zaharia, Andrew Chen, Aaron Davidson, Ali Ghodsi, Sue~Ann Hong, Andy
  Konwinski, Siddharth Murching, Tomas Nykodym, Paul Ogilvie, Mani Parkhe,
  et~al.
\newblock Accelerating the machine learning lifecycle with {MLflow}.
\newblock \emph{IEEE Data Engineering Bulletin}, 41(4):39--45, 2018.

\end{thebibliography}

\end{document}